\section{Planning：LLM Agent 的核心挑战}

\subsection{什么是 Planning}

\begin{importantbox}{Planning 的核心含义}
Planning 是对任务的分解——不管是隐式还是显式的，都要：
\begin{enumerate}
    \item 识别用户的目标
    \item 分解任务为可执行的子步骤
    \item Think Ahead：提前想多步
    \item Think Multiple Stages：考虑多个阶段
\end{enumerate}
\end{importantbox}

引入 Planning 的原因在于：语言模型在解决复杂的、长周期的、Multi-Hop 的问题时存在明显缺陷。通过显式或隐式的 Planning 机制，可以 Leverage 模型能力，最大化其在复杂任务上的 Performance。

\subsection{Planning 方法分类 Survey}

根据 2024 年初的一篇经典 Survey（关于 LLM 做 Planning 的综述），可以将方法分为以下几大类：

\subsubsection{分解（Decomposition）}

\begin{itemize}
    \item \textbf{CoT（Chain of Thought）}：``Let's think step by step''，让模型一步步求解问题，内部涉及对问题的自由分解。
    \item \textbf{ReAct}：不断和环境交互——Thought $\rightarrow$ Action $\rightarrow$ Observation 循环，在交互过程中完成任务分解。
    \item \textbf{Plan-and-Solve}：显式地形成一个 Plan，然后沿着 Plan 去 Solve。
    \item \textbf{PoT（Program of Thought）}：将思考过程转化为可执行的程序。
\end{itemize}

\subsubsection{选择（Selection）}

\begin{itemize}
    \item \textbf{ToT（Tree of Thoughts）}：用树结构管理搜索过程，从一个节点出发扩展多个候选方案，基于评分选择最有潜力的节点进一步探索。
    \item \textbf{MCTS（Monte Carlo Tree Search）}：基于 UCB 值综合考虑探索和利用，选择节点进行拓展。关键在于 \textbf{Simulation}（对节点进行打分）和 \textbf{Selection}（选择最优节点）。
\end{itemize}

\begin{warningbox}{状态爆炸问题}
在 ToT 和 MCTS 中，从每个节点出发语言模型可以产生 multiple 的 proposal，这会导致指数级的状态爆炸。因此必须引入有效的评分机制（Score），对节点做剪枝（Prune），只对有潜力的节点进一步探索。
\end{warningbox}

\subsubsection{外部求解器（External Solver）}

如 PDDL Planner，用于 Robotics 等场景中的 Planning 任务。语言模型直接产生 Plan 时往往不 work，引入外部符号求解器可以大幅提升可靠性。

\subsubsection{反思（Reflection \& Critique）}

通过 Self-Refine 机制，模型对自身输出进行反思和改进，在记忆（Memory）层面也可以引入 Reflection 来形成更高层次的结论。

\begin{knowledgebox}{Generative Agents 中的 Memory Reflection}
在 Generative Agents 中，如果只记录浅层事实（"他在读论文"），价值有限。通过触发 Reflection——综合多次观察，可以形成更高层结论（"这是一个对科研非常有热情的人"），这才是有深度的记忆。
\end{knowledgebox}

\subsection{Planning 方法的个人分类}

作者按自己的维度将 Planning 方法分为五类：

\begin{table}[H]
\centering
\begin{tabular}{ll}
\toprule
\textbf{类别} & \textbf{代表方法} \\
\midrule
Reasoning Model（隐式） & DeepSeek R1, OpenAI O 系列 \\
Plan-based（显式） & ReAct, Plan-and-Solve \\
Search-based & ToT, LLM + MCTS \\
Test-Time Scaling & Majority Voting, DeepConf \\
Neuro-Symbolic & LM+P, Alpha Geometry, PDDL \\
\bottomrule
\end{tabular}
\caption{Planning 方法分类}
\end{table}

\subsection{本章小结}

LLM 在面对复杂长周期问题时，直接生成 Plan 容易出错。可以通过分解（CoT/ReAct）、搜索（ToT/MCTS）、外部求解器（PDDL）、反思（Self-Critique）等多种策略来提升 Planning 质量。


